\documentclass[main.tex]{subfiles}


\begin{document}

%from pagenumber to up
% \phantomsection 
% \hypertarget{toc}{}


 
   


\section{Как найти ответ к задаче}

В начале изучения физики приходится решать задачи, но так как опыта нет, то для упражнения нужно научиться просто находить ответ к классическим задачам. На примере следующей задачи покажем, как это можно делать на основе общих соображений.

\medskip

\noindent\textbf{Задача.} Трактор за первые 5~мин проехал 600~м. Какой путь он пройдет за 0,5~ч, двигаясь с той же скоростью? 

\medskip

Мы знаем, что для решения нам нужно начать с формулы. Смотрим в пособие по физике и видим, что путь~$S$ равен
\begin{equation}
\label{find_e1}
S=vt,
\end{equation}
где $v,t$~--- скорость тела и время его движения.

К чему это? Можно по ней что-то найти. Мы видим связанные 5~мин и 600~м, и их подстановка в формулу дает
$$
600=v\cdot5\quad \Rightarrow\quad v=120~\text{м/мин}.
$$

Так по частям можно решать любую задачу. Что дальше? Ведь нужно найти <<какой-то путь>>\ldots\

Можно понять что найти-то нужно путь <<второй>> (600~м~--- это <<первый>> путь).

Но любой путь при равномерном движении находится по формуле~\eqref{find_e1}. Для второго пути есть время, не хватает для него скорости. Но скорость-то будет такая же! Находим второй путь:
$$
S_2=120\cdot30=3600~\text{м}
$$
(тут перевели 0,5~ч в 30~мин, чтоб <<подходило>>).

\smallskip

Здесь вкратце показан простейший пример творческих рассуждений при решении задачи по физике. Развитие умения творчески рассуждать, а не решать по алгоритмам и~т.\,п., поможет впоследствии перейти к задачам вплоть до олимпиадного уровня. 

{\small
\begin{quote}
    \ldots математика есть весьма подсобный предмет для физики, она очень помогает, но это, во всяком случае, не есть что-то необходимое.

    (П.\,Капица. «Как следует изучать физику».)
\end{quote}
\begin{quote}
    \ldots Вы правильно пришли к выводу, что надо меньше думать об основах. Главное, чем надо овладеть,~--- это техникой работы, а понимание тонкостей само придет потом.

    (Л.\,Ландау. Из писем.)
\end{quote}
\begin{quote}
    \ldots путь к ответу~--- это индивидуальный и увлекательный научный поиск. И этот творческий процесс нельзя заменить изучением рецептов решения задач.

    (О.\,Я.\,Савченко. <<Задачи по физике>>.)
\end{quote}
    
}



\end{document}
